\documentclass{article}
\usepackage[utf8]{inputenc}
\usepackage[T1]{fontenc}
\usepackage{geometry}
\geometry{
    a4paper,
    left=2.5cm,
    right=2.5cm,
    top=2.5cm,
    bottom=2.5cm
}
\usepackage{graphicx}
\usepackage{float}
\usepackage{hyperref}
\usepackage[english]{babel}
\usepackage{amsmath}
\usepackage{amssymb}
\usepackage{array}
\usepackage{booktabs}
\usepackage{caption}
\usepackage{subcaption}
\usepackage{enumitem}
\usepackage{fancyhdr}
\usepackage{titlesec}
\usepackage{url}
\usepackage{textcomp}
\usepackage{xcolor}
\usepackage{listings}
\usepackage{multirow}
\usepackage{longtable}
\usepackage{tabularx}
\usepackage{times}

% Set up page style
\pagestyle{fancy}
\fancyhf{}
\rfoot{\thepage}
\renewcommand{\headrulewidth}{0pt}

% Title formatting
\titleformat{\chapter}[display]{\normalfont\bfseries}{\chaptertitle\thechapter}{0pt}{\Huge}
\titlespacing*{\chapter}{0pt}{50pt}{40pt}

% Document info
\title{Template Document}
\author{Academic Support Agent}
\date{\today}

\begin{document}
\maketitle

%  ,  \maketitle 
%  cls 
\maketitle

\newpage
%  
\makeabstract

 
 
 Binaural
Modeling 

......

\myspace{1}%  
\noindent { \blackti  :} HRTF 

%  
\makeEnabstract

Virtual auditory technology is also called auralization. It is brought forward as a new concept with the development of acoustic simulation techniques in recent years and can be implemented by establishing the physical or mathematical models of a sound field to achieve sound effects simulation. If we consider the binaural effect, it can be called binaural virtual auditory.

......

\myspace{1}
\noindent \textbf{ KEY WORDS }virtual auditory, HRTF, neural network

%  
\makecontent
\iffalse
    \makefigcontent
    \maketabcontent
\fi
%  
\maketext

\section{ }
\subsection{ }
\subsubsection{ }

 Auralization \cite{Xu_2023_ICCV} \cite{li2021learning} Audio rendering Acoustical simulation 
\makespace
\section{ }
\subsection{ }
\subsubsection{ }
\begin{equation}
    \int_a^b f(x)\mathrm{d}x=F(b)-F(a)
\end{equation}

\subsubsection{ }
 (\ref{q1}) 
\begin{equation}
    \label{q1}
    \begin{aligned}
        dx      & =v_{x}dt  \\
        dy      & =v_{y}dt  \\
        x_{t+1} & =dx+x_{t} \\
        y_{t+1} & =dx+y_{t}
    \end{aligned}
\end{equation}

\subsubsection{ }
 (\ref{q1}) 
\begin{equation}
    \left\{
    \begin{aligned}
        100(t-kT_{2}) & , & t\in (kT_{2}, kT_{2}+0.2 )    \\
        20            & , & t\in (kT_{2}+0.2,kT_{2}+2.2)  \\
        -100t+240     & , & t\in( kT_{2}+2.2, kT_{2}+2.4) \\
        0             & , & t\in (kT_{2}+2.4,(k+1)T_{2})
    \end{aligned}
    \right.
\end{equation}

\subsection{ }
\subsubsection{ }
 \ref{table1} 

\begin{table}[H]
    \fontsize{10.5bp}{1.25}
    \caption{ }
    % \setlength{\abovecaptionskip}{0.2cm}  % 
    % \setlength{\belowcaptionskip}{0.2cm} % 
    \centering\label{table1}
    % \vspace{-0.5em}	
    \begin{tabular}{|c|c|c|c|}
        \hline % \toprule[1.5pt]
                & A  & B  & C  \\ \hline
         /dB    & 0.86  & 1.02  & 0.69  \\ \hline
         /s & 25    & 25    & 27    \\ \hline
    \end{tabular}
\end{table}

\subsubsection{ }
 \ref{table3}
\begin{table}[H]
    \fontsize{10.5bp}{1.25}
    \caption{ }
    % \setlength{\abovecaptionskip}{0.2cm}  % 
    % \setlength{\belowcaptionskip}{0.2cm} % 
    \centering\label{table3}
    % \vspace{-0.5em}	
    \begin{tabular}{cccc}
        \toprule[1.5pt]
                & A  & B  & C  \\
        \toprule[1.5pt]
         /dB    & 0.86  & 1.02  & 0.69  \\
         /s & 25    & 25    & 27    \\
        \toprule[1.5pt]
    \end{tabular}
\end{table}

\subsubsection{ }
 \ref{table2}
\begin{table}[H]
    \caption{ }
    % \setlength{\abovecaptionskip}{0.2cm}  % 
    % \setlength{\belowcaptionskip}{-0.2cm} % 
    \fontsize{10.5bp}{1.25}
    \centering\label{table2}
    % \vspace{-0.5em}	
    \begin{tabular}{|c|c|c|c|}
        \hline % \toprule[1.5pt]
        \makecell*[c]{ Parameter Group     \\ Condition Selection } &\makecell*[c]{ Basic Ways \\of Hatching  } 	& \makecell*[c]{Calculated Average\\ Alapsed Time} & \makecell*[c]{Calculated Average\\ Alapsed Time}  \\
        \hline
        \makecell*[c]{ Parameter group (1) \\ \\Parameter group (2)\\} &\makecell*[c]{ Zigzag   Hatch\\Contour Hatch\\Zigzag   Hatch\\Contour Hatch}& \makecell*[c]{468.940\\374.923 \\885.792\\1947.77}   & \makecell*[c]{Zigzag 1.888 \\ Contour  5.195   }   \\
        \hline
        \makecell*[c]{ Parameter group (1) \\ \\ Parameter group (2)\\}  &\makecell*[c]{ Zigzag   Hatch\\Contour Hatch\\Zigzag   Hatch\\Contour Hatch}&\makecell*[c]{ 545.080\\356.847\\1068.275\\1692.098} & \makecell*[c]{Zigzag 1.960 \\ Contour 4.742  }  \\
        \hline
    \end{tabular}
    % \vspace{-1.5em}
\end{table}

\subsection{ }
\subsubsection{ }
 \ref{P1} 
\begin{figure}[!ht]
    \centering
    %\vspace{-0.30cm} % `
    %\setlength{\abovecaptionskip}{0.0cm}   % `
    %\setlength{\belowcaptionskip}{0.0cm} % 
    \includegraphics[width=4cm]{./picture/test.png}
    \caption{ }\label{P1}
\end{figure}
\subsubsection{ }
 \ref{P2} \ref{P2.1} \ref{P2.2} 
\begin{figure}[h!]
    \centering
    \begin{subfigure}[b]{0.25\linewidth}
        \includegraphics[width=\linewidth]{picture/test.png}
        \caption{ 1}\label{P2.1}
    \end{subfigure}
    \begin{subfigure}[b]{0.25\linewidth}
        \includegraphics[width=\linewidth]{picture/test.png}
        \caption{ 2}\label{P2.2}
    \end{subfigure}
    \caption{ }
    \label{P2}
\end{figure}
\makespace
\section{ Test}
\subsection{ Test}
\subsubsection{ Test}
%  
\makespace
\section*{ }
\begingroup  %  thebibliography 
\renewcommand{\section}[2]{}
\addcontentsline{toc}{section}{ }
% npu 

\addtolength{\itemsep}{-0.6em} %  
\begin{spacing}{1.0}           % 
    
\end{spacing}
%  

%  
\makespace % 
\section*{ }
\begin{center}
    { \blackti \fontsize{16.0600pt}{1.25}  \,  }
\end{center}
\addcontentsline{toc}{section}{ \texorpdfstring{ \, }{}  }
\myspace{1}
 

%  
\section*{ }
\makespace
\begin{center}
    { \blackti \fontsize{16.0600pt}{1.25} }
\end{center}
\addcontentsline{toc}{section}{ }
\myspace{1}
 

%  
\makespace
\section*{ }
\begin{center}
    { \blackti \fontsize{16.0600pt}{1.25}  \,  }
\end{center}
\addcontentsline{toc}{section}{ \texorpdfstring{ \, }{}  }
\myspace{1}
\end{document}